\documentclass[11pt,a4paper]{article}

\usepackage[a4paper,margin=1in]{geometry}
\usepackage{amsmath}
\usepackage{booktabs}
\usepackage{array}
\usepackage{enumitem}
\usepackage{tabularx}
\usepackage{fancyhdr}
\usepackage[T1]{fontenc}
\usepackage{lmodern}
\usepackage{verbatim}

\newcommand{\universityname}{UNIVERSITY OF DERBY}
\newcommand{\college}{College of Science and Engineering}
\newcommand{\examtitle}{Timed In-Class Test Worked Solutions}
\newcommand{\moduletitle}{OPERATIONAL RESEARCH}
\newcommand{\modulecode}{5CM522}
\newcommand{\examdate}{MAY 2026}

\pagestyle{fancy}
\fancyhf{}
\fancyfoot[C]{\thepage}
\renewcommand{\headrulewidth}{0pt}

\setlength{\parindent}{0pt}
\setlength{\parskip}{0.45em}
\setlist[itemize]{leftmargin=1.1cm, itemsep=0.25em, topsep=0.25em}
\setlist[enumerate,1]{label=(\alph*), leftmargin=1.1cm, itemsep=0.4em, topsep=0.3em}

\begin{document}

\begin{center}
{\large \textbf{\universityname}}\\
{\college}\\[0.4em]
{\Large \textbf{\examtitle}}\\
{\large \textbf{\moduletitle} -- \textbf{\modulecode}}\\[0.3em]
\textbf{DATE:} \examdate
\end{center}

These solutions match the current live paper.

Source paper:
\begin{itemize}
\item \texttt{paper/in-class-test-2026-05.tex}
\end{itemize}

Runnable verification helpers are in:
\begin{itemize}
\item \texttt{solutions/in\_class\_test\_code.py}
\item \texttt{solutions/worked\_solution\_snippets.py}
\end{itemize}

To verify the whole solution pack quickly:
\begin{verbatim}
python solutions/worked_solution_snippets.py
\end{verbatim}

\textbf{1. Linear Programming}

\textbf{a) Model formulation}

Let:
\[
x = \text{number of Standard panels}, \qquad y = \text{number of Premium panels}
\]

Model:
\[
\max Z = 42x + 36y
\]
subject to
\[
4x + 2y \le 84
\]
\[
2x + 3y \le 60
\]
\[
y \le 18
\]
\[
x, y \ge 0
\]

For marking, either assumption should be accepted if stated clearly:
\begin{itemize}
\item continuous non-negative decision variables, giving the LP solution below
\item whole-number decision variables, since panels are indivisible in practice
\end{itemize}

\textbf{b) Optimal weekly production plan}

At the optimum, printing and finishing are both binding:
\[
4x + 2y = 84, \qquad 2x + 3y = 60
\]

Solving gives:
\[
x = 16.5, \qquad y = 9
\]

Maximum weekly contribution:
\[
Z = 42(16.5) + 36(9) = 1017
\]

\textbf{c) Binding constraints and unused capacity}

Binding/slack:
\begin{itemize}
\item Printing slack = 0
\item Finishing slack = 0
\item Premium market slack = $18 - 9 = 9$
\end{itemize}

\textbf{d) Value of extra capacity}

Extra-capacity interpretation:
\begin{itemize}
\item With 6 extra printing hours, contribution rises to 1057.5
\item Improvement = 40.5, so value per printing hour = 6.75
\item With 6 extra finishing hours, contribution rises to 1062
\item Improvement = 45, so value per finishing hour = 7.50
\end{itemize}

So finishing is more valuable at the current optimum, and the maximum worthwhile price is:
\[
\pounds 7.50 \text{ per finishing hour, for up to the first 6 hours}
\]

\textbf{e) Premium-share requirement}

Premium-share requirement:
\[
y \ge 0.3(x + y)
\]
which rearranges to
\[
0.7y \ge 0.3x \qquad \Longleftrightarrow \qquad y \ge \frac{3x}{7}
\]

At the original optimum:
\[
\frac{y}{x+y} = \frac{9}{16.5 + 9} = \frac{9}{25.5} \approx 35.29\%
\]

So the original optimum is still feasible, and the optimal plan remains unchanged.

If a student explicitly assumes integer production quantities, the corresponding answers are:
\begin{itemize}
\item Optimal whole-number plan: $x = 17$, $y = 8$
\item Contribution: $Z = 42(17) + 36(8) = 1002$
\item Slack values:
\begin{itemize}
\item Printing slack = 0
\item Finishing slack = $60 - (2(17) + 3(8)) = 2$
\item Premium market slack = $18 - 8 = 10$
\end{itemize}
\item With 6 extra printing hours, the best integer contribution is 1050
\item With 6 extra finishing hours, the best integer contribution is 1062
\item So finishing is still more valuable over this 6-hour comparison. The average gain is \pounds 8 per printing hour versus \pounds 10 per finishing hour, but this should not be treated as a constant shadow price in the integer model.
\item The 30\% Premium requirement is still satisfied because $\frac{8}{17+8} = 32\%$, so the integer optimum also remains unchanged.
\end{itemize}

Verification snippet:
\begin{verbatim}
from worked_solution_snippets import q1_results
print(q1_results())
\end{verbatim}

\textbf{2. Binary Integer Programming}

\textbf{a) Model formulation}

Let $x_A, \ldots, x_F$ be binary variables where $x_s = 1$ if site $s$ is opened.

Baseline model:
\[
\min Z = 16x_A + 11x_B + 15x_C + 14x_D + 13x_E + 10x_F
\]

Coverage constraints:
\[
x_A + x_B \ge 1
\]
\[
x_A + x_C + x_E \ge 1
\]
\[
x_B + x_C + x_F \ge 1
\]
\[
x_A + x_D + x_F \ge 1
\]
\[
x_C + x_D + x_E \ge 1
\]
\[
x_D + x_E + x_F \ge 1
\]
\[
x_s \in \{0,1\}
\]

\textbf{b) Baseline optimal solution}

Baseline optimum:
\[
\text{Open } B, E, F
\]
Minimum weekly cost:
\[
34
\]

\textbf{c) Redundant coverage in the baseline solution}

Redundancy in the baseline solution:
\begin{itemize}
\item District 3 is covered by $B$ and $F$
\item District 6 is covered by $E$ and $F$
\end{itemize}

Operational value:
\begin{itemize}
\item This redundancy gives resilience if one selected site is temporarily unavailable.
\end{itemize}

\textbf{d) Revised model and impact of restrictions}

Revised model adds:
\[
x_D + x_E + x_F \ge 2
\]
\[
x_B + x_F \le 1
\]

Revised optimum:
\[
\text{Open } B, D, E
\]
Minimum weekly cost:
\[
38
\]

\textbf{e) One realistic extension}

One realistic extension:
\begin{itemize}
\item Add site capacity limits or maximum response times.
\end{itemize}

Extra data needed:
\begin{itemize}
\item expected demand by district
\item travel times from each site to each district
\item maximum launch capacity at each site
\end{itemize}

Verification snippet:
\begin{verbatim}
from worked_solution_snippets import q2_results
print(q2_results())
\end{verbatim}

\textbf{3. Transhipment}

\textbf{a) Transhipment model formulation}

Let:
\[
x_{ph} = \text{flow from depot } p \text{ to hub } h
\]
\[
y_{hc} = \text{flow from hub } h \text{ to clinic } c
\]

Objective:
\begin{itemize}
\item Minimise total inbound and outbound transport cost.
\end{itemize}

Constraints:
\begin{itemize}
\item Supply constraints at depots
\item Demand constraints at clinics
\item Flow balance at each hub
\item Non-negativity
\end{itemize}

\textbf{b) Baseline optimal shipment plan}

Baseline optimal solution:
\begin{itemize}
\item $P_1 \to H_1 = 100$
\item $P_2 \to H_2 = 120$
\item $P_3 \to H_2 = 40$
\item $H_1 \to C_1 = 70$
\item $H_1 \to C_4 = 30$
\item $H_2 \to C_2 = 60$
\item $H_2 \to C_3 = 80$
\item $H_2 \to C_4 = 20$
\end{itemize}

Baseline total cost:
\[
1950
\]

Unused supply:
\[
40 \text{ units at } P_3
\]

\textbf{c) Disrupted scenario}

Disrupted model adds:
\[
\sum_c y_{H_1 c} \le 130
\]
\[
y_{H_1 C_2} \ge 30
\]

Revised optimal solution:
\begin{itemize}
\item $P_1 \to H_1 = 100$
\item $P_2 \to H_2 = 120$
\item $P_3 \to H_2 = 40$
\item $H_1 \to C_1 = 20$
\item $H_1 \to C_2 = 30$
\item $H_1 \to C_4 = 50$
\item $H_2 \to C_1 = 50$
\item $H_2 \to C_2 = 30$
\item $H_2 \to C_3 = 80$
\end{itemize}

Revised total cost:
\[
2040
\]

Interpretation:
\begin{itemize}
\item Forcing at least 30 units for clinic $C_2$ through $H_1$ changes the network pattern and increases cost.
\end{itemize}

\textbf{d) Emergency direct route from $P_3$ to $C_4$}

Emergency direct route extension:
\begin{itemize}
\item Add direct decision variable $d$ for $P_3 \to C_4$
\item Bounds: $0 \le d \le 20$
\item Add $9d$ to the objective
\item The 50 comes from clinic $C_4$'s original demand, so replace the $C_4$ demand constraint with $y_{H_1 C_4} + y_{H_2 C_4} + d = 50$
\item Reduce available $P_3$ hub shipments accordingly, so $x_{P_3 H_1} + x_{P_3 H_2} + d \le 80$
\end{itemize}

Optimal direct-shipping decision:
\[
d = 20
\]

Revised non-zero shipments:
\begin{itemize}
\item $P_1 \to H_1 = 100$
\item $P_2 \to H_2 = 120$
\item $P_3 \to H_2 = 20$
\item $P_3 \to C_4 = 20$
\item $H_1 \to C_1 = 70$
\item $H_1 \to C_2 = 30$
\item $H_2 \to C_2 = 30$
\item $H_2 \to C_3 = 80$
\item $H_2 \to C_4 = 30$
\end{itemize}

Revised total cost with emergency direct route:
\[
2020
\]

So the direct route is worthwhile because it reduces cost by 20. It is cheaper than sending those units via $H_2$, since $P_3 \to H_2 \to C_4$ costs $4 + 6 = 10$ per unit while the emergency route costs 9 per unit.

Verification snippet:
\begin{verbatim}
from worked_solution_snippets import q3_results
print(q3_results())
\end{verbatim}

\textbf{4. Simulation}

\textbf{a) Main DES elements}

Main DES elements:
\begin{itemize}
\item Entities: students
\item Resource: technicians
\item Events: arrivals, service start, service completion
\item Stochastic inputs: interarrival times and service times
\item Measures: number served, average wait, maximum wait, percentage waiting less than 10 minutes
\end{itemize}

\textbf{b) Baseline SimPy run with 2 technicians}

Baseline one-run results with 2 technicians:
\begin{itemize}
\item Students served = 71
\item Average wait = 5.2208 minutes
\item Maximum wait = 46.4285 minutes
\item Percentage under 10 minutes = 87.3239\%
\end{itemize}

\textbf{c) Busy-period staffing comparison}

Busy-period results at 13 students per hour:

\begin{center}
\begin{tabular}{lrrrr}
\toprule
Technicians & Served & Avg wait & Max wait & \% under 10 min \\
\midrule
2 & 110 & 43.8859 & 170.5482 & 30.9091 \\
3 & 105 & 4.0800 & 28.1272 & 83.8095 \\
4 & 105 & 0.1727 & 3.2712 & 100.0000 \\
\bottomrule
\end{tabular}
\end{center}

Recommendation from part (c):
\begin{itemize}
\item 3 technicians is the minimum staffing level that meets the primary target of average wait below 8 minutes.
\item It does not meet the secondary target of at least 85\% under 10 minutes.
\end{itemize}

\textbf{d) Mixed service-time extension}

Mixed-service extension:
\begin{itemize}
\item Standard queries: 75\%, mean service 6 minutes
\item Complex queries: 25\%, mean service 14 minutes
\item Keep arrivals at 13 per hour
\item For each arrival, sample whether the case is standard or complex
\item Then sample service time from the corresponding exponential distribution
\end{itemize}

One run with 3 technicians gives:
\begin{itemize}
\item Students served = 84
\item Complex cases = 29
\item Average wait = 2.1404 minutes
\item Maximum wait = 14.5512 minutes
\item Percentage under 10 minutes = 91.6667\%
\end{itemize}

So the 3-technician recommendation still looks reasonable on this run.

\textbf{e) Why one simulation run is not enough}

Why one run is not enough:
\begin{itemize}
\item Output depends on random variation.
\item One run can be unusually good or bad.
\item Repeated runs allow a mean estimate and confidence interval.
\end{itemize}

Verification snippet:
\begin{verbatim}
from worked_solution_snippets import q4_results
print(q4_results())
\end{verbatim}

\textbf{5. Forecasting}

\textbf{a) Centred moving averages and trend}

Using periods 1--12, the centred 4-quarter moving averages can be calculated only for periods 3--10:

\begin{center}
\begin{tabular}{rrrr}
\toprule
Period & Actual & MA4 & CMA \\
\midrule
2  & 505 & 482.5000 &  \\
3  & 455 & 490.0000 & 486.2500 \\
4  & 590 & 501.2500 & 495.6250 \\
5  & 410 & 512.5000 & 506.8750 \\
6  & 550 & 526.2500 & 519.3750 \\
7  & 500 & 535.0000 & 530.6250 \\
8  & 645 & 547.5000 & 541.2500 \\
9  & 445 & 557.5000 & 552.5000 \\
10 & 600 & 572.5000 & 565.0000 \\
\bottomrule
\end{tabular}
\end{center}

Fitting a straight line to the CMA values for periods 3--10 gives the estimated trend:
\[
\text{Trend} = 451.1756 + 11.3095t
\]

\textbf{b) Seasonal indices}

Estimated seasonal indices:

Additive:
\[
Q1 = -101.8750, \quad Q2 = 33.1250, \quad Q3 = -30.6250, \quad Q4 = 99.3750
\]

Multiplicative:
\[
Q1 = 0.8076, \quad Q2 = 1.0611, \quad Q3 = 0.9396, \quad Q4 = 1.1917
\]

\textbf{c) Validation-set comparison using MSE}

Validation-set accuracy on periods 13--16 uses:
\[
\mathrm{MSE} = \frac{1}{n}\sum_{t=1}^{n} e_t^2
\]
where
\[
e_t = \text{Actual}_t - \text{Forecast}_t
\]

Results:
\[
\mathrm{MSE}_{\text{additive}} = 558.4018
\]
\[
\mathrm{MSE}_{\text{multiplicative}} = 184.4082
\]

So the better model is multiplicative.

\textbf{d) Forecasts for 2023 and recommendation}

Forecasts for 2023 using the multiplicative model:

\begin{center}
\begin{tabular}{lcr}
\toprule
Period & Quarter & Forecast \\
\midrule
17 & Q1 & 519.6556 \\
18 & Q2 & 694.7350 \\
19 & Q3 & 625.7974 \\
20 & Q4 & 807.2476 \\
\bottomrule
\end{tabular}
\end{center}

Recommendation:
\begin{itemize}
\item Use the multiplicative model.
\item It gives the lower validation MSE.
\item The seasonal swing appears to rise with the level of the series, which supports a multiplicative interpretation.
\end{itemize}

Observation to check:
\begin{itemize}
\item Period 16 with demand 775 is the largest observation and worth checking before operational use.
\end{itemize}

Verification snippet:
\begin{verbatim}
from worked_solution_snippets import q5_results
print(q5_results())
\end{verbatim}

\textbf{6. Scheduling}

\textbf{Part A: Johnson's Rule}

\textbf{a) Johnson's Rule sequence}

Processing times:
\[
J_1 = (8,4), \quad J_2 = (3,9), \quad J_3 = (6,5), \quad J_4 = (7,2), \quad J_5 = (4,6)
\]

Johnson steps:
\begin{enumerate}
\item Smallest time is 2 on $J_4$ on Machine 2, so place $J_4$ last.
\item Next smallest is 3 on $J_2$ on Machine 1, so place $J_2$ first.
\item Next smallest is 4 on $J_5$ on Machine 1, so place $J_5$ next.
\item Next smallest remaining is 4 on $J_1$ on Machine 2, so place $J_1$ in the latest remaining slot.
\item $J_3$ fills the remaining middle slot.
\end{enumerate}

Optimal Johnson sequence:
\[
J_2, J_5, J_3, J_1, J_4
\]

\textbf{b) Gantt chart and minimum makespan}

Gantt-style timing:

\begin{center}
\begin{tabular}{lrrrr}
\toprule
Job & M1 start & M1 finish & M2 start & M2 finish \\
\midrule
$J_2$ & 0  & 3  & 3  & 12 \\
$J_5$ & 3  & 7  & 12 & 18 \\
$J_3$ & 7  & 13 & 18 & 23 \\
$J_1$ & 13 & 21 & 23 & 27 \\
$J_4$ & 21 & 28 & 28 & 30 \\
\bottomrule
\end{tabular}
\end{center}

Minimum makespan:
\[
30
\]

\textbf{c) Comparison with the original order}

Original-order schedule $J_1, J_2, J_3, J_4, J_5$ gives makespan:
\[
34
\]

Improvement:
\[
4 \text{ time units}
\]

\textbf{Part B: Job Shop Solution}

\textbf{d) Job shop model, schedule and makespan}

Use an optimisation model such as OR-Tools CP-SAT with:
\begin{itemize}
\item Objective: minimise makespan
\item Decision variables: start times, end times, or interval tasks for each operation
\item Precedence constraints: each task must follow the previous task in the same job
\item Machine-capacity constraints: no overlap on each machine
\end{itemize}

One optimal baseline schedule is:
\begin{center}
\begin{tabular}{lrlrr}
\toprule
Job & Task & Machine & Start & End \\
\midrule
$A$ & 1 & $M_1$ & 0  & 3  \\
$B$ & 1 & $M_2$ & 0  & 4  \\
$C$ & 1 & $M_3$ & 0  & 3  \\
$A$ & 2 & $M_3$ & 3  & 8  \\
$B$ & 2 & $M_1$ & 4  & 6  \\
$C$ & 2 & $M_2$ & 4  & 8  \\
$D$ & 1 & $M_1$ & 6  & 8  \\
$B$ & 3 & $M_3$ & 8  & 14 \\
$C$ & 3 & $M_1$ & 8  & 11 \\
$D$ & 2 & $M_2$ & 8  & 14 \\
$A$ & 3 & $M_2$ & 14 & 16 \\
$D$ & 3 & $M_3$ & 14 & 18 \\
\bottomrule
\end{tabular}
\end{center}

So the minimum makespan is:
\[
18
\]

Short interpretation:
\begin{itemize}
\item Precedence is enforced because each job's next operation cannot start before the previous one finishes.
\item Non-overlap is enforced because each machine can process at most one operation at a time.
\end{itemize}

\textbf{e) Revised feasibility with deadline and machine downtime}

Part (e) revised model:
\begin{itemize}
\item Add due-date constraint: completion time of job $D$ must be $\le 15$
\item Add machine downtime on $M_2$ from time 6 to 9, for example by inserting a blocking interval
\end{itemize}

Re-solving the model gives:
\[
\text{infeasible}
\]

Operational interpretation:
\begin{itemize}
\item The new restrictions remove too much flexibility from the schedule.
\item Job $D$ needs 6 uninterrupted time units on $M_2$ after first visiting $M_1$.
\item Because $M_2$ is unavailable from time 6 to 9, job $D$ cannot both complete its $M_2$ operation and still finish the full job by time 15.
\item A solver reporting infeasibility means there is no schedule that satisfies all constraints simultaneously.
\end{itemize}

Verification snippet:
\begin{verbatim}
from worked_solution_snippets import q6_results
print(q6_results())
\end{verbatim}

\end{document}
